\documentclass{amsart}
% For dealing with references we use the comment environment
\usepackage{verbatim}
\newenvironment{reference}{\comment}{\endcomment}
%\newenvironment{reference}{}{}
\newenvironment{slogan}{\comment}{\endcomment}
\newenvironment{history}{\comment}{\endcomment}

% For commutative diagrams we use Xy-pic
\usepackage[all]{xy}

% We use 2cell for 2-commutative diagrams.
\xyoption{2cell}
\UseAllTwocells

% We use multicol for the list of chapters between chapters
\usepackage{multicol}

% This is generally recommended for better output
\usepackage{lmodern}
\usepackage[T1]{fontenc}

% For cross-file-references

% Package for hypertext links:
\usepackage{hyperref}

% For any local file, say "hello.tex" you want to link to please

% Theorem environments.
%
\theoremstyle{plain}
\newtheorem{theorem}[subsection]{Theorem}
\newtheorem{proposition}[subsection]{Proposition}
\newtheorem{lemma}[subsection]{Lemma}

\theoremstyle{definition}
\newtheorem{definition}[subsection]{Definition}
\newtheorem{example}[subsection]{Example}
\newtheorem{exercise}[subsection]{Exercise}
\newtheorem{situation}[subsection]{Situation}

\theoremstyle{remark}
\newtheorem{remark}[subsection]{Remark}
\newtheorem{remarks}[subsection]{Remarks}

\numberwithin{equation}{subsection}

% Macros
%
\def\lim{\mathop{\mathrm{lim}}\nolimits}
\def\colim{\mathop{\mathrm{colim}}\nolimits}
\def\Spec{\mathop{\mathrm{Spec}}}
\def\Hom{\mathop{\mathrm{Hom}}\nolimits}
\def\Ext{\mathop{\mathrm{Ext}}\nolimits}
\def\SheafHom{\mathop{\mathcal{H}\!\mathit{om}}\nolimits}
\def\SheafExt{\mathop{\mathcal{E}\!\mathit{xt}}\nolimits}
\def\Sch{\mathit{Sch}}
\def\Mor{\mathop{\mathrm{Mor}}\nolimits}
\def\Ob{\mathop{\mathrm{Ob}}\nolimits}
\def\Sh{\mathop{\mathit{Sh}}\nolimits}
\def\NL{\mathop{N\!L}\nolimits}
\def\CH{\mathop{\mathrm{CH}}\nolimits}
\def\proetale{{pro\text{-}\acute{e}tale}}
\def\etale{{\acute{e}tale}}
\def\QCoh{\mathit{QCoh}}
\def\Ker{\mathop{\mathrm{Ker}}}
\def\Im{\mathop{\mathrm{Im}}}
\def\Coker{\mathop{\mathrm{Coker}}}
\def\Coim{\mathop{\mathrm{Coim}}}

% Boxtimes
%
\DeclareMathSymbol{\boxtimes}{\mathbin}{AMSa}{"02}

%
% Macros for moduli stacks/spaces
%
\def\QCohstack{\mathcal{QC}\!\mathit{oh}}
\def\Cohstack{\mathcal{C}\!\mathit{oh}}
\def\Spacesstack{\mathcal{S}\!\mathit{paces}}
\def\Quotfunctor{\mathrm{Quot}}
\def\Hilbfunctor{\mathrm{Hilb}}
\def\Curvesstack{\mathcal{C}\!\mathit{urves}}
\def\Polarizedstack{\mathcal{P}\!\mathit{olarized}}
\def\Complexesstack{\mathcal{C}\!\mathit{omplexes}}
% \Pic is the operator that assigns to X its picard group, usage \Pic(X)
% \Picardstack_{X/B} denotes the Picard stack of X over B
% \Picardfunctor_{X/B} denotes the Picard functor of X over B
\def\Pic{\mathop{\mathrm{Pic}}\nolimits}
\def\Picardstack{\mathcal{P}\!\mathit{ic}}
\def\Picardfunctor{\mathrm{Pic}}
\def\Deformationcategory{\mathcal{D}\!\mathit{ef}}

\setlength{\emergencystretch}{3em}
\begin{document}
\title[Stacks Project, Tag 0EL1]{499 --- Lefschetz cohomology comparison through the punctured affine cone}
\maketitle
\noindent Copyright (C) 2005--2025 Johan de Jong. Stacks Project authors. Source: \href{https://stacks.math.columbia.edu/tag/0EL1}{Tag 0EL1}.

\noindent Editorial difficulty note (not author text). Difficulty H3: The argument transfers the depth inequality to the cone through the two possible fibre-point types of a one-dimensional smooth torsor. Local cohomology annihilation and completion yield a comparison on the cone, and compatible graded decompositions recover the required proper-scheme cohomology comparison..

\noindent Editorial provenance note (not author text). History: excerpt from revision \texttt{a0444\allowbreak 6e57e\allowbreak c1fbc\allowbreak 252a8\allowbreak 71afc\allowbreak ec775\allowbreak 2fb28\allowbreak 07b14}, algebraization.tex, lines 9020--9152. Reference presentation was adapted; original-author context and core support blocks were retained or added. The mathematical wording is unchanged. Licensed under GNU FDL 1.2 or later; no invariant sections or cover texts. See ../licenses/COPYING. Context blocks reproduce author definitions and supporting results; they are not additional counted problems.

\begin{lemma}
\label{more-morphisms-lemma-apply-proj-spec}
Let $R$ be a ring. Let $P$ be a proper scheme over $R$ and let
$\mathcal{L}$ be an ample invertible $\mathcal{O}_P$-module.
Set $A = \bigoplus_{m \geq 0} \Gamma(P, \mathcal{L}^{\otimes m})$.
Then $P = \text{Proj}(A)$ and diagram (\href{https://stacks.math.columbia.edu/tag/0EKG}{equation proj and spec (Tag 0EKG)})
becomes the diagram
$$
\xymatrix{
\underline{\Spec}_P \left(
\bigoplus\nolimits_{m \in \mathbf{Z}} \mathcal{L}^{\otimes m}
\right)
\ar[r] \ar@{=}[d] &
L =
\underline{\Spec}_P \left(
\bigoplus\nolimits_{m \geq 0} \mathcal{L}^{\otimes m}
\right) \ar[d]^\sigma \ar[r]_-\pi & P \ar[d] \\
U \ar[r] & X \ar[r] & Z
}
$$
having the properties explained above.
\end{lemma}

\medskip\noindent
We continue to assume $P$ is quasi-compact.
Let $\mathcal{F}$ be a quasi-coherent $\mathcal{O}_P$-module.
Let us set $\mathcal{F}_U = \pi^*\mathcal{F}|_U$. Then we have
\begin{equation}
\label{more-morphisms-equation-cohomology-torsor}
R\Gamma(U, \mathcal{F}_U) =
\bigoplus\nolimits_{m \in \mathbf{Z}}
R\Gamma(P, \mathcal{F} \otimes_{\mathcal{O}_P} \mathcal{O}_P(m))
\end{equation}
Moreover, this direct sum decomposition is functorial in $\mathcal{F}$
and the induced $A$-module structure on the right is the same as the
$A$-module structure on the left coming from $U \subset X$.
To prove the formula, since $\pi|_U$ is affine and
$(\pi|_U)_*\mathcal{O}_U = \bigoplus_{m \in \mathbf{Z}} \mathcal{O}_P(m)$
we get
\begin{align*}
R(\pi|_U)_*\mathcal{F}_U
& =
(\pi|_U)_*\mathcal{F}_U \\
& =
(\pi|_U)_*(\pi|_U)^*\mathcal{F} \\
& =
\mathcal{F} \otimes_{\mathcal{O}_P} 
\bigoplus\nolimits_{m \in \mathbf{Z}} \mathcal{O}_P(m) \\
& =
\bigoplus\nolimits_{m \in \mathbf{Z}}
\mathcal{F} \otimes_{\mathcal{O}_P} \mathcal{O}_P(m)
\end{align*}
By Leray we find that
$R\Gamma(U, \mathcal{F}_U) =
R\Gamma(P, R(\pi|_U)_*\mathcal{F}_U)$, see
Cohomology, Lemma \href{https://stacks.math.columbia.edu/tag/01F4}{lemma apply Leray (Tag 01F4)}.
The proof is finished because taking cohomology commutes with
direct sums in this case, see Derived Categories of Schemes, Lemma
\href{https://stacks.math.columbia.edu/tag/08DZ}{lemma quasi coherence pushforward direct sums (Tag 08DZ)}.
This is where we use that $P$ is quasi-compact; $P$ is separated
by Constructions, Lemma \href{https://stacks.math.columbia.edu/tag/01MC}{lemma proj separated (Tag 01MC)}.

\begin{definition}
\label{divisors-definition-zero-scheme-s}
Let $X$ be a scheme. Let $\mathcal{L}$ be an invertible sheaf.
Let $s \in \Gamma(X, \mathcal{L})$ be a global section.
The {\it zero scheme} of $s$ is the closed subscheme $Z(s) \subset X$
defined by the quasi-coherent sheaf of ideals
$\mathcal{I} \subset \mathcal{O}_X$ which is the image of the
map $s : \mathcal{L}^{\otimes -1} \to \mathcal{O}_X$.
\end{definition}

\medskip\noindent
Let $k$ be a field. Let $P$ be a proper scheme over $k$.
Let $\mathcal{L}$ be an ample invertible $\mathcal{O}_P$-module.
Let $s \in \Gamma(P, \mathcal{L})$ be a section\footnote{We do not
require $s$ to be a regular section. Correspondingly, $Q$ is only
a locally principal closed subscheme of $P$ and not necessarily an effective
Cartier divisor.} and let
$Q = Z(s)$ be the zero scheme, see
Divisors, Definition \ref{divisors-definition-zero-scheme-s}.
For all $n \geq 1$ we denote $Q_n = Z(s^n)$
the $n$th infinitesimal neighbourhood of $Q$.
If $\mathcal{F}$ is a coherent $\mathcal{O}_P$-module, then we denote
$\mathcal{F}_n = \mathcal{F}|_{Q_n}$ the restriction, i.e.,
the pullback of $\mathcal{F}$ by the closed immersion
$Q_n \to P$.


\medskip\noindent
Let $R$ be a ring. Let $A$ be a graded $R$-algebra, see
Algebra, Section \href{https://stacks.math.columbia.edu/tag/00JL}{section graded (Tag 00JL)}.
For $m \geq 0$ we denote $A_{\geq m} = \bigoplus_{d \geq m} A_d$.
Consider the graded ring
$$
B = \bigoplus\nolimits_{d \geq 0} A_{\geq d}
$$
For $d' \geq d$ and $a \in A_{d'}$ let us denote $a^{(d)} \in B$
the element in $B_d$ corresponding to $a$.
Let us denote $\sigma : A \to B$ and $\psi : A \to B$
the two obvious ring maps: if $a \in A_d$, then
$\sigma(a) = a^{(0)}$ and $\psi(a) = a^{(d)}$.
Then $\psi$ is a graded ring map and $\sigma$ turns $B$
into a graded algebra over $A$. There is also a
surjective graded ring map $\tau : B \to A$
which for $d' \geq d$ and $a \in A_{d'}$ sends
$a^{(d)}$ to $0$ if $d' > d$ and to $a$ if $d' = d$.

\medskip\noindent
Affine schemes and spectra.
We set $X = \Spec(A)$. The irrelevant ideal $A_+$ cuts out a closed subscheme
$Z = V(A_+) = \Spec(A/A_+) = \Spec(A_0)$. Set $U = X \setminus Z$.
$$
U \longrightarrow X \longrightarrow Z
$$
Projective schemes and Proj. Set $P = \text{Proj}(A)$. We may
and do view $P$ as a scheme over $\Spec(A_0) = Z$.
Set $L = \text{Proj}(B)$. We may and do view $L$ as a scheme
over $\Spec(B_0) = \Spec(A) = X$; observe that the identification
of $B_0$ with $A$ is given by $\sigma$.
The surjection $\tau$ defines a closed immersion $0 : P \to L$.
Since $A \xrightarrow{\sigma} B \to A$ is equal to the map $A \to A_0 \to A$
we conclude that
$$
\xymatrix{
P \ar[d] \ar[r]_0 & L \ar[d] \\
Z \ar[r] & X
}
$$
is commutative.


\begin{proposition}
\label{proposition-lefschetz}
In the situation above assume there exists an integer $\sigma$ such that
for all points $p \in P \setminus Q$ we have
$$
\text{depth}(\mathcal{F}_p) + \dim(\overline{\{p\}}) > \sigma
$$
Then the map
$$
H^i(P, \mathcal{F}) \longrightarrow \lim H^i(Q_n, \mathcal{F}_n)
$$
is an isomorphism for $0 \leq i < \sigma$.
\end{proposition}

\begin{proof}
We will use More on Morphisms, Lemma \ref{more-morphisms-lemma-apply-proj-spec}
and we will use the notation used and results found
More on Morphisms, Section \href{https://stacks.math.columbia.edu/tag/0EKF}{section proj spec (Tag 0EKF)}
without further mention; this proof will not make sense without
at least understanding the statement of the lemma.
Observe that in our case
$A = \bigoplus_{m \geq 0} \Gamma(P, \mathcal{L}^{\otimes m})$
is a finite type $k$-algebra
all of whose graded parts are finite dimensional $k$-vector spaces, see
Cohomology of Schemes, Lemma \href{https://stacks.math.columbia.edu/tag/0B5T}{lemma coherent proper ample (Tag 0B5T)}.

\medskip\noindent
We may and do think of $s$ as an element $f \in A_1 \subset A$, i.e.,
a homogeneous element of degree $1$ of $A$. Denote
$Y = V(f) \subset X$ the closed subscheme defined by $f$.
Then $U \cap Y = (\pi|_U)^{-1}(Q)$ scheme theoretically.
Recall the notation
$\mathcal{F}_U = \pi^*\mathcal{F}|_U = (\pi|_U)^*\mathcal{F}$.
This is a coherent $\mathcal{O}_U$-module.
Choose a finite $A$-module $M$ such that
$\mathcal{F}_U = \widetilde{M}|_U$
(for existence see Local Cohomology, Lemma
\href{https://stacks.math.columbia.edu/tag/0BK0}{lemma finiteness pushforwards and H1 local (Tag 0BK0)}).
We claim that $H^i_Z(M)$ is annihilated by
a power of $f$ for $i \leq \sigma + 1$.

\medskip\noindent
To prove the claim we will apply
Local Cohomology, Proposition \href{https://stacks.math.columbia.edu/tag/0EFC}{proposition annihilator (Tag 0EFC)}.
Translating into geometry we see that it suffices to prove
for $u \in U$, $u \not \in Y$ and $z \in \overline{\{u\}} \cap Z$
that
$$
\text{depth}(\mathcal{F}_{U, u}) +
\dim(\mathcal{O}_{\overline{\{u\}}, z}) > \sigma + 1
$$
This requires only a small amount of thought.

\medskip\noindent
Observe that $Z = \Spec(A_0)$ is a finite set of closed points of $X$ because
$A_0$ is a finite dimensional $k$-algebra.
(The reader who would like $Z$ to be a singleton can replace the
finite $k$-algebra $A_0$ by $k$; it won't affect anything else in the proof.)

\medskip\noindent
The morphism $\pi : L \to P$ and its restriction $\pi|_U : U \to P$
are smooth of relative dimension $1$.
Let $u \in U$, $u \not \in Y$ and $z \in \overline{\{u\}} \cap Z$.
Let $p = \pi(u) \in P \setminus Q$ be its image.
Then either $u$ is a generic
point of the fibre of $\pi$ over $p$ or a closed point of the fibre.
If $u$ is a generic point of the fibre, then
$\text{depth}(\mathcal{F}_{U, u}) = \text{depth}(\mathcal{F}_p)$
and
$\dim(\overline{\{u\}}) = \dim(\overline{\{p\}}) + 1$.
If $u$ is a closed point of the fibre, then
$\text{depth}(\mathcal{F}_{U, u}) = \text{depth}(\mathcal{F}_p) + 1$
and
$\dim(\overline{\{u\}}) = \dim(\overline{\{p\}})$.
In both cases we have
$\dim(\overline{\{u\}}) = \dim(\mathcal{O}_{\overline{\{u\}}, z})$
because every point of $Z$ is closed. Thus the desired
inequality follows from the assumption in the statement of
the lemma.

\medskip\noindent
Let $A'$ be the $f$-adic completion of $A$. So $A \to A'$ is flat by
Algebra, Lemma \href{https://stacks.math.columbia.edu/tag/00MB}{lemma completion flat (Tag 00MB)}.
Denote $U' \subset X' = \Spec(A')$ the inverse image of
$U$ and similarly for $Y'$ and $Z'$. Let $\mathcal{F}'$
on $U'$ be the pullback of $\mathcal{F}_U$ and let
$M' = M \otimes_A A'$.
By flat base change for local cohomology
(Local Cohomology, Lemma \href{https://stacks.math.columbia.edu/tag/0EF5}{lemma torsion change rings (Tag 0EF5)})
we have
$$
H^i_{Z'}(M') = H^i_Z(M) \otimes_A A'
$$
and we find that for $i \leq \sigma + 1$
these are annihilated by a power of $f$. Consider the diagram
$$
\xymatrix{
& H^i(U, \mathcal{F}_U) \ar[ld] \ar[d] \ar[r] &
\lim H^i(U, \mathcal{F}_U/f^n\mathcal{F}_U) \ar@{=}[d] \\
H^i(U, \mathcal{F}_U) \otimes_A A' \ar@{=}[r] &
H^i(U', \mathcal{F}') \ar[r] &
\lim H^i(U', \mathcal{F}'/f^n\mathcal{F}')
}
$$
The lower horizontal arrow is an isomorphism for $i < \sigma$ by
Lemma \href{https://stacks.math.columbia.edu/tag/0EKM}{lemma alternative higher (Tag 0EKM)} and the torsion
property we just proved. The horizontal equal sign is flat base change
(Cohomology of Schemes, Lemma \href{https://stacks.math.columbia.edu/tag/02KH}{lemma flat base change cohomology (Tag 02KH)})
and the vertical equal sign is because $U \cap Y$ and $U' \cap Y'$
as well as their $n$th infinitesimal neighbourhoods
are mapped isomorphically onto each other (as we are
completing with respect to $f$).

\medskip\noindent
Applying More on Morphisms, Equation
(\ref{more-morphisms-equation-cohomology-torsor})
we have compatible direct sum decompositions
$$
\lim H^i(U, \mathcal{F}_U/f^n\mathcal{F}_U) =
\lim
\left(
\bigoplus\nolimits_{m \in \mathbf{Z}}
H^i(Q_n, \mathcal{F}_n \otimes \mathcal{L}^{\otimes m})
\right)
$$
and
$$
H^i(U, \mathcal{F}_U) =
\bigoplus\nolimits_{m \in \mathbf{Z}}
H^i(P, \mathcal{F} \otimes \mathcal{L}^{\otimes m})
$$
Thus we conclude by Algebra, Lemma \href{https://stacks.math.columbia.edu/tag/0EKD}{lemma daniel litt (Tag 0EKD)}.
\end{proof}
\end{document}
